% from ega3-6-fr.tex
\oldpage[III]{143}

des applications canoniques $P_\bullet^\lambda\to P_\bullet$,
$Q_\bullet^\mu\to Q_\bullet$, forment un système inductif d'homomorphismes,
d'où un homomorphisme canonique (6.3.6.1), et plus généralement un
homomorphisme canonique
\[
  \varinjlim_{\lambda,\mu}''E(P_\bullet^\lambda,Q_\bullet^\mu)
  \longrightarrow ''E(P_\bullet,Q_\bullet)
\]
dont (6.3.6.1) est l'homomorphisme des aboutissements. En outre, la suite
spectrale $''E(P_\bullet,Q_\bullet)$ est régulière (6.3.2), et il en est de
même de la suite spectrale
$\varinjlim_{\lambda,\mu}''E(P_\bullet^\lambda,Q_\bullet^\mu)$, comme il
résulte des définitions (0, 11.1.7) et de la démonstration de (0, 11.3.3);
pour démontrer que (6.3.6.1) est bijectif, il suffit donc (0, 11.1.5) de
prouver que l'homomorphisme
\[
  \varinjlim_{\lambda,\mu}''E(P_\bullet^\lambda,Q_\bullet^\mu)
  \longrightarrow ''E(P_\bullet,Q_\bullet)
  \tag{6.3.6.2}\label{III.6.3.6.2.fr}
\]
est bijectif pour les termes $E^2$. Comme le foncteur $H_\bullet$ commute à
la limite inductive des complexes de modules, on est finalement ramené à
prouver que pour deux systèmes inductifs filtrants $(F^\lambda)$,
$(G^\mu)$ de $A$-modules, l'homomorphisme canonique
\[
  \varinjlim_{\lambda,\mu}\operatorname{Tor}_\bullet^A(F^\lambda,G^\mu)
  \longrightarrow
  \operatorname{Tor}_\bullet^A
  \left(\varinjlim_\lambda F^\lambda,
        \varinjlim_\mu G^\mu\right)
\]
est bijectif. Pour cela, considérons pour chaque $F^\lambda$ la résolution
libre canonique
\[
  L_\bullet^\lambda:\quad
  \cdots\longrightarrow L_{i+1}^\lambda
  \longrightarrow L_i^\lambda\longrightarrow\cdots
  \longrightarrow L_1^\lambda\longrightarrow L_0^\lambda\longrightarrow0
\]
où $L_0^\lambda$ est le $A$-module des combinaisons linéaires formelles
d'éléments de $F^\lambda$ et $L_{i+1}^\lambda$ le $A$-module des combinaisons
linéaires formelles d'éléments de
$\operatorname{Ker}(L_i^\lambda\to L_{i-1}^\lambda)$; on vérifie
immédiatement que les $L_\bullet^\lambda$ forment un système inductif de
complexes, et si l'on pose
$F=\varinjlim_\lambda F^\lambda$, $L_i=\varinjlim_\lambda L_i^\lambda$, les
$L_i$ forment une résolution $L_\bullet$ de $F$, le foncteur
$\varinjlim$ étant exact; en outre, les $L_i$, limites inductives de
$A$-modules libres, sont plats (0\textsubscript{I}, 6.1.2). On considère de
même pour chaque $\mu$ la résolution libre canonique $M_\bullet^\mu$ de
$G^\mu$, et $M_\bullet=\varinjlim_\mu M_\bullet^\mu$ est une résolution plate
de $G=\varinjlim_\mu G^\mu$. On a alors
$\operatorname{Tor}_\bullet^A(\varinjlim_\lambda F^\lambda,
\varinjlim_\mu G^\mu)=H_\bullet(L_\bullet\otimes_A M_\bullet)$ en vertu de
(6.3.5) et (6.3.4); mais
\[
  H_\bullet(L_\bullet\otimes_A M_\bullet)
  =\varinjlim_{\lambda,\mu}
   H_\bullet(L_\bullet^\lambda\otimes_A M_\bullet^\mu)
\]
puisque $H_\bullet$ commute aux limites inductives de complexes de modules;
comme $H_\bullet(L_\bullet^\lambda\otimes_A M_\bullet^\mu)
=\operatorname{Tor}_\bullet^A(F^\lambda,G^\mu)$, cela termine la
démonstration.

Lorsqu'on suppose qu'il existe $i_0$ tel que
$P_i^\lambda=Q_i^\mu=0$ pour $i<i_0$ quels que soient $\lambda$ et $\mu$, on
démontre de la même manière que l'homomorphisme canonique
\[
  \varinjlim_{\lambda,\mu}'E(P_\bullet^\lambda,Q_\bullet^\mu)
  \longrightarrow'E(P_\bullet,Q_\bullet)
  \tag{6.3.6.3}\label{III.6.3.6.3.fr}
\]
est bijectif.
\end{proof}

\begin{proposition}[6.3.7]
Supposons $P_\bullet$ et $Q_\bullet$ limités inférieurement. Si le complexe
$P_\bullet$ est formé de $A$-modules plats, on a un $A$-isomorphisme
canonique de $\partial$-foncteurs en $Q_\bullet$
\[
  \mathbf{Tor}_\bullet^A(P_\bullet,Q_\bullet)
  \xrightarrow{\sim}H_\bullet(P_\bullet\otimes_A Q_\bullet).
  \tag{6.3.7.1}\label{III.6.3.7.1.fr}
\]
\end{proposition}

\begin{proof}
En effet, la suite spectrale (6.3.2.1) est birégulière et dégénérée, et
l'existence de l'isomorphisme (6.3.7.1) résulte de (0, 11.1.6). En outre, en
calculant l'hypertor à partir d'une résolution projective de Cartan-Eilenberg
de $P_\bullet$ (6.3.4), on voit aussitôt que l'isomorphisme ainsi défini est
un isomorphisme de $\partial$-foncteurs en $Q_\bullet$.
\end{proof}
\begin{env}[6.3.8]
